% !TeX spellcheck = de_DE

Ein essenzieller Part von iterativer KI gestützten Verbesserung von Leichter Sprache + Texten ist die automatische Textvereinfachung (Text Simplification). Sie zielt darauf ab, den Inhalt und die Struktur eines Textes so zu verändern, dass dieser einfacher zu verstehen und lesen ist, während die ursprüngliche Bedeutung beibehalten wird \cite[1. Introduction]{alva-manchego_data-driven_2020}. \\
Dabei werden hauptsächlich vier Arten von Texttransformation durchgeführt:
\begin{itemize}[label=\textbullet, nosep]
	\item \textbf{Ersetzung} schwieriger Wörter durch einfachere Synonyme
	\item \textbf{Umordnung} der Satzstruktur
	\item \textbf{Löschen} von unwichtigen Informationen
	\item \textbf{Satzspaltung} langer Sätze in kürzere
\end{itemize}

Die Forschung unterscheidet hierbei zur Umsetzung zwischen verschiedenen Methodischen Ansätzen:
\begin{table}[h!]
	\centering
	\begin{tabularx}{\textwidth}{|X|X|}
		\hline
		\textbf{Methodischer Ansatz:} & \textbf{Erklärung} \\
		\hline
		Statistische maschinelle Übersetzung (SMT) & Betrachtet Vereinfachung als Übersetzung von der „komplexen“ in die „einfache“ Quellsprache. Nutzt statistische Wahrscheinlichkeitsmodelle, die auf parallelen Korpora trainiert wurden \cite[S. 154]{alva-manchego_data-driven_2020}. \\
		\hline
		Phrasenbasierte Ansätze (PBSMT) & Eine Weiterentwicklung von SMT, bei der nicht mehr einzelne Wörter, sondern ganze Wortfolgen (\textit{Phrasen}) als fundamentale Einheiten übersetzt und transformiert werden. Dies erlaubt eine bessere Beibehaltung des lokalen Kontextes \cite[S. 154]{alva-manchego_data-driven_2020}. \\
		\hline
	\end{tabularx}
	\caption{Datengesteuerte Ansätze (Data-Driven Approaches)}
\end{table}